% Section 6.2: Inner, outer and contraction products | Part: 1 | Depends on: sec:clifford | Status: done | Original lines: 1411-1462
%!TEX root = ../../../main.tex
\subsection{Inner, outer and contraction products}

\begin{definition}[Products induced by the geometric product]\label{def:products}
For homogeneous multivectors $A_r\in\GA^r$ and $B_s\in\GA^s$:
\begin{align*}
 A_r\wedge B_s&:=\gr{A_rB_s}{r+s}&&\text{(outer product)},\\
 A_r\lc B_s&:=\gr{A_rB_s}{s-r}&&\text{(left contraction; $=0$ if $r>s$)},\\
 A_r*B_s&:=\gr{A_rB_s}{0}&&\text{(scalar product)},
\end{align*}
extended bilinearly to all multivectors. The squared norm (Euclidean case) is $\norm A^2=\gr{A\rev A}{0}\ge0$.
\end{definition}

For a vector $a$ and an $r$-vector $A_r$ one has the explicit formulas
\begin{equation}\label{eq:aA}
 aA_r=a\lc A_r+a\wedge A_r,\qquad
 a\lc A_r=\tfrac12\bigl(aA_r-(-1)^rA_ra\bigr),\quad
 a\wedge A_r=\tfrac12\bigl(aA_r+(-1)^rA_ra\bigr),
\end{equation}
of grades $r-1$ and $r+1$ respectively (for $r=1$ this is \eqref{eq:abdecomp}).

\begin{proposition}[Properties]\label{prop:contr}
For $a\in V$ and multivectors $A,B,C$:
\begin{enumerate}
  \item (explicit formula) $a\lc(b_1\wedge\cdots\wedge b_k)=\sum_{i=1}^k(-1)^{i-1}(a\cdot b_i)\,b_1\wedge\cdots\wedge\widehat{b_i}\wedge\cdots\wedge b_k$.
  \item (antiderivation) $a\lc(A_s\wedge B)=(a\lc A_s)\wedge B+(-1)^sA_s\wedge(a\lc B)$.
  \item (associativity-type identity) $A\lc(B\lc C)=(A\wedge B)\lc C$.
  \item $(A\wedge B)\wedge C=A\wedge(B\wedge C)$, and $A_r\wedge B_s=(-1)^{rs}B_s\wedge A_r$.
\end{enumerate}
\end{proposition}
\begin{proof}[Proof idea]
(1) and (2) follow by induction from \eqref{eq:aA} (for $k=2$: $a\lc(b\wedge c)=(a\cdot b)c-(a\cdot c)b$). (3) is checked first for $A=a$ a vector, using (2), and then holds by induction on the grade of $A$. (4) is the statement that the wedge defined here agrees with the exterior product of \S\ref{sec:ext} (Corollary~\ref{cor:grades}).
A complete proof of (3), on basis blades, is Identity~\ref{id:contr3} in the Glossary.% REWRITTEN: glossary pointer
\end{proof}
% AUTHOR-NOTE (resolved): (3) is only sketched here, and Theorem thm:dual depends on it; the full proof is now
% Identity id:contr3 in Glossary Part C (docs/AUTHOR_NOTES.md #6).

\begin{example}
$a\lc(b\wedge c)=(a\cdot b)\,c-(a\cdot c)\,b$ is a vector lying in the plane $b\wedge c$ and perpendicular to $a$: indeed $a\cdot\bigl((a\cdot b)c-(a\cdot c)b\bigr)=0$.
\end{example}

\begin{remark}[Commutator product]
For bivectors $A,B$ one has $AB=A\lc B+\tfrac12(AB-BA)+A\wedge B$ with the three terms of grades $0,2,4$. The antisymmetric middle part, $A\times B:=\tfrac12(AB-BA)$ is again a bivector; it makes $\GA^2$ a Lie algebra isomorphic to $\mathfrak{so}(p,q)$.
\end{remark}

\begin{intuition}
\begin{itemize}
  \item \textbf{Wedge = ``span''.} $A\wedge B$ builds the higher-dimensional element containing both (it is $0$ if they share a direction). It is the \emph{join}.
  \item \textbf{Contraction = ``the part of $B$ perpendicular to $A$, measured by $A$''.} $a\lc B$ removes the direction $a$ from the blade $B$ and returns what remains inside $B$, perpendicular to $a$, with the scale of $|a|$. It lowers grade by one.
  \item \textbf{Geometric product = both at once.} For a vector $a$ and blade $B$, $aB=a\lc B+a\wedge B$ splits into the \emph{lower}-grade part (the part of $a$ lying inside $B$) and the \emph{higher}-grade part (the part of $a$ sticking out of $B$). Because the splitting is a sum of different grades, no information is lost, and that is why $a$ can be divided out again.
\end{itemize}
\end{intuition}
